

Une masse $m_0 = 1.00 ~ \si{\kilo\gram}$ d’eau sous forme de glace à \SI{0.00}{\celsius} est transformée totalement en eau vapeur à \SI{100}{\celsius}, sous la pression atmosphérique.

\Q Proposer un schéma de la transformation, étape par étape.

\Q Calculer les variations $\Delta H$ et $\Delta U$ du corps pur.

\textbf{Données :}
\begin{align*}
& \Delta_{fus} h (\ce{H2O}, 0 ~ \si{\celsius}) = 334 ~ \si{\kilo\joule\per\kilo\gram} ;  \Delta_{vap} h (\ce{H2O}, 100 ~ \si{\celsius}) = 2260 ~ \si{\kilo\joule\per\kilo\gram} ; \\
& \rho(\ce{H2O}, s) = 9.18 \cdot 10^2 ~ \si{\kilo\gram\per\meter\cubed} ; p(\ce{H2O}, l) = 1.00 x 10^3 ~ \si{\kilo\gram\per\meter\cubed} ; \\
& c_e = 4.18 ~ \si{\kilo\joule\per\kilo\gram}.
\end{align*}

